\subsection{The exact minimal-polynomial target in the cubic example}

P-F36 distinguishes the algebraic coefficient from a finite Laurent
character. The worked example itself permits a more precise positive
comparison after retaining its quadratic extension. The positivity
recorded below for this one coefficient was displayed computationally
in GCT VIII; the elementary identities make its scope explicit.

\paragraph{The source's normalization and target.}
Use $A,D,P,x,c$ from P-F36, with the exact displayed discriminant
$D=(q^8+1)^2q^{24}A$. GCT VIII, Section~5.3, pp.~31--32,
defines $v_1^1$ by replacing $Q_\sigma$ with $P_\sigma$ in
$u_1^1$, without changing the scalar coefficients. Thus the coefficient
$c$ of $P_{1212121}$ in $v_1^1$ is precisely (P-F36c).
Section~5.4.1, p.~48, explains that these convenient coefficients must
be rescaled by suitable powers of the positive Laurent polynomial
$\widehat f_p$ displayed in GCT VII, Section~7.1.7, p.~28.
It studies their traces and norms, equivalently their quadratic
minimal-polynomial coefficients. Figure~21, p.~55, and
Figures~22--23, pp.~56--57, display the resulting coefficient vectors.

Set $F=q^{11}\widehat f_p$. The monomial factor merely centers its
exponents. To specify this normalization without a convention about
the ordering of the displayed source terms, write
\[
 F(q)=f_0+\sum_{j=1}^{16}f_j(q^j+q^{-j}),\qquad
 (f_0,\ldots,f_{16})=
 (16,16,16,16,17,18,18,18,16,14,12,10,8,6,4,2,1).
 \tag{P-F39a}
\]
In particular $F(1)=400$ and $F$ has positive coefficients.
No unimodality of $F$ itself is needed or asserted.
For a finite vector $v=(v_0,\ldots,v_d)$, abbreviate
\[
 [v](q)=v_0+\sum_{j=1}^d v_j(q^j+q^{-j}).
\]

\begin{proposition}[P-F39: positive minimal-polynomial coefficients and the remaining comparison]
\label{prop:nonstandardtracenorm}
Let $a=Fc\in\mathbb C(q,x)$, with $x^2=D$.
\begin{enumerate}
\item The minimal polynomial of $a$ over $\mathbb C(q)$ is
\[
                   Z^2+T_1(q)Z+N_1(q),
 \qquad T_1=FT,\quad N_1=F^2N,
 \tag{P-F39b}
\]
where
\[
 \begin{split}
 T(q)&=22-2(q^2+q^{-2})+6(q^4+q^{-4})
       +2(q^6+q^{-6})\\
     &\hspace{15mm}+3(q^8+q^{-8})+2(q^{10}+q^{-10}),\\
 N(q)&=138-14(q^2+q^{-2})+70(q^4+q^{-4})
      +20(q^6+q^{-6})\\
     &\quad+31(q^8+q^{-8})+26(q^{10}+q^{-10})
      +6(q^{12}+q^{-12})\\
     &\quad+9(q^{14}+q^{-14})+3(q^{16}+q^{-16})
      +3(q^{18}+q^{-18})+(q^{20}+q^{-20}).
 \end{split}
 \tag{P-F39c}
\]
Both $T_1$ and $N_1$ are symmetric Laurent polynomials with positive
integer coefficients, nonincreasing from the middle towards either
end. Their values at one are respectively $17600$ and $71680000$.
Each is the Tate-degree character of a finite direct sum of
Tate-twisted cohomologies of projective spaces, with their usual
Lefschetz operators.
\item There is an explicit finite \emph{\'{e}tale} quadratic cover
of a nonempty open subset of the $q$-line whose rank-two pushforward
of the structure sheaf has a multiplication endomorphism with
characteristic polynomial (P-F39b). Nevertheless this endomorphism
cannot be identified with a horizontal endomorphism of a finite
complex local system on that open subset, or after any connected
finite dominant cover. Its trace is nonconstant.
\end{enumerate}
The cohomological character realizations in the first assertion
identify the two rational coefficients separately. They do not
identify a nonstandard block, multiplication map, or theta-lift
projector with those cohomology objects.
\end{proposition}

\begin{proof}
The conjugate of $c$ is $c'=-(P-x)/(2q^{20})$. Therefore
\[
       -(c+c')=P/q^{20}=T,
 \qquad cc'=(P^2-D)/(4q^{40})=N.
 \tag{P-F39d}
\]
Expanding gives exactly (P-F39c). Since $F$ is a nonzero element
of the rational-function field, P-F36 implies that $a=Fc$ is
not rational. It has degree two, so its characteristic polynomial
$(Z-Fc)(Z-Fc')$ is also its minimal polynomial. This proves
(P-F39b).

For completeness, the entire finite coefficient calculation is
given here. Multiplication of (P-F39a) by (P-F39c) gives
$T_1=[\alpha]$ and $N_1=[\beta]$, where
\[
\begin{split}
\alpha={}&(708,708,694,680,673,666,642,618,565,512,464,416,365,314,\\
&262,210,170,130,106,82,62,42,30,18,11,4,2),
\end{split}
 \tag{P-F39e}
\]
and
\[
\begin{split}
\beta={}&(2212002,2205476,2186350,2155076,2112378,2058980,1995675,\\
&1923256,1842975,1756084,1663849,1567536,1468102,1366504,\\
&1263855,1161268,1059834,960644,864459,772040,684050,601152,\\
&523831,452572,387494,328716,276117,229576,188839,153652,\\
&123592,98236,77108,59732,45640,34364,25488,18596,13346,9396,\\
&6493,4384,2894,1848,1142,672,377,196,95,40,15,4,1).
\end{split}
 \tag{P-F39f}
\]
These identities can be checked coefficient by coefficient using
the ordinary convolution rule. More explicitly, put $f_{-j}=f_j$
for $|j|\le16$, and $f_j=0$ otherwise. Put
$t_{\pm 2j}=t_{2j}$ with $(t_0,t_2,\ldots,t_{10})=(22,-2,6,2,3,2)$,
and set the other $t_j$ equal to zero. Define $n_j$ in the same
way from the coefficients in (P-F39c). Then the displayed entries
are exactly
\[
       \alpha_j=\sum_i f_i t_{j-i},\qquad
       \beta_j=\sum_{i,k} f_i f_k n_{j-i-k}.
 \tag{P-F39g}
\]
The support bounds make every sum finite. These formulas,
together with the explicit lists, verify the polynomial identities
without interpolation or numerical approximation. Every entry
of both lists is positive and each list is nonincreasing. Their
values at one can either be summed from the lists or obtained
from $F(1)=400$, $T(1)=44$, and $N(1)=448$.

Here is an actual cohomological realization of the resulting
characters. For any nonincreasing nonnegative integral vector
$v=(v_0,\ldots,v_d)$, put $v_{d+1}=0$ and
$m_j=v_j-v_{j+1}$. Telescoping gives
\[
 [v](q)=\sum_{j=0}^d m_j q^{-j}(1+q+\cdots+q^{2j}).
 \tag{P-F39h}
\]
The coefficient of $q^k$ on the right is
$\sum_{j\ge|k|}(v_j-v_{j+1})=v_{|k|}$, which proves
the identity for every exponent. The cohomology of
$\mathbb P^{2j}$ consists of one Tate class in each degree
$0,2,\ldots,4j$. Its twist by $(j)$ has Tate-degree
character $q^{-j}(1+q+\cdots+q^{2j})$, where
$\mathbb Q_\ell(-i)$ is recorded as $q^i$.
This is the projective-space case of Milne's projective-bundle
formula, \emph{Lectures on \'{E}tale Cohomology}, Theorem~23.2,
printed pp.~139--140. The hyperplane class gives the usual
Lefschetz chain. Thus
\[
       \bigoplus_{j=0}^d
       H^*(\mathbb P^{2j},\mathbb Q_\ell)(j)^{\oplus m_j}
 \tag{P-F39i}
\]
realizes $[v]$ with nonnegative multiplicities and even Tate
cohomology. Apply this to $v=\alpha$ and $v=\beta$.
This assertion concerns the specified direct sums and their
Tate grading, and does not require all summands to have one
common cohomological degree or weight.

To retain the quadratic datum geometrically, let
\[
 R=\mathbb C[q,q^{-1},D^{-1},F^{-1}],\qquad
 S=R[x]/(x^2-D),\qquad U=\operatorname{Spec}R.
\]
The extension is free of rank two with basis $1,x$.
Because $D$ is invertible, so is $x$, and the derivative
$2x$ of its defining monic polynomial is invertible.
Hence $\operatorname{Spec}S\to U$ is finite \'{e}tale.
It is connected: P-F36 proves irreducibility over the
fraction field, and $S$ embeds in that quadratic field
since it is free over $R$.
Multiplication by $a=-F(P+x)/(2q^{20})$ on $S$ is represented by
\[
 -\frac{F}{2q^{20}}
   \begin{pmatrix}P&D\\1&P\end{pmatrix}.
 \tag{P-F39j}
\]
Taking its determinant polynomial gives exactly (P-F39b).
This is an endomorphism of an algebraic vector bundle,
with the quadratic cover and its involution retained.

Finally an endomorphism of a complex local system has,
in local flat frames, locally constant matrices. Its trace
is therefore locally constant, and is constant on a connected
base. Equivalently a horizontal endomorphism of the associated
flat bundle has constant characteristic-polynomial coefficients.
Here the trace is $-T_1(q)$, which is nonconstant: for example
its coefficient at $q^{26}$ is $2$, by (P-F39e).
On any connected finite dominant cover of $U$, at least one
irreducible component dominates $U$. The inclusion of function
fields on that component is injective and fixes the constants,
so a nonconstant rational function on $U$ remains nonconstant
there. In particular its pullback is not constant on the cover.
Consequently the multiplication endomorphism cannot become
a horizontal local-system endomorphism after such a cover.
This proves the second assertion.
\end{proof}

\paragraph{What a sign or purity argument would still have to identify.}
Before the prescribed rescaling, the rational minimal-polynomial
coefficients $T$ and $N$ already satisfy the source's positive-at-one
condition, but have negative Laurent coefficients. After one power
of the exact scaling factor, the two coefficients in (P-F39b)
have positive unimodal characters, and (P-F39i) is a concrete
Tate--Lefschetz model for each. Their construction here is a
coefficient-level realization; it supplies no correspondence
with the nonstandard multiplication, cell modules, or automorphic
projectors. Likewise the finite cover in (P-F39j) supplies the
algebraic eigenvalue datum, while its varying multiplication
operator is not a morphism of the cover's local system of
constant coefficients.

This locates a specific missing comparison rather than ruling
out local systems, algebraic coefficient fields, or theta lifts
in general. A construction may treat $q$ as a formal grading
variable rather than as a coordinate on a base; the horizontal
obstruction applies only to the latter literal interpretation.
No asymptotic bound in tensor degree, or conjecture for all
nonstandard coefficients, follows from this fixed calculation.
The formal saturation and almost-positivity targets are
GCT VIII Conjectures~3.5--3.6, pp.~24--25; the worked
monomial-basis calculation is not silently identified with a
coefficient in a different ambient canonical basis.

Primary sources are
\url{https://arxiv.org/pdf/0709.0749}, Section~7.1.7, p.~28;
\url{https://arxiv.org/pdf/0709.0751}, Section~5.3, pp.~31--32,
Figure~6, p.~38, Section~5.4.1, p.~48, and Figures~21--23,
pp.~55--57; and
\url{https://www.jmilne.org/math/CourseNotes/LEC.pdf},
Theorem~23.2, pp.~139--140.
